% ===================================================================
%  TABLICA Nr 3 — Dzieciństwo i młodość.
%  Meme regime que les tableaux 1 et 2.
%
%  LA SCENE EST UN BLOC A ELLE, ET SON ORDINAL SUIT LE NOM :
%  « Scena pierwsza », non « Pierwsza scena ». C'est la place que
%  html.py a du apprendre pour cette colonne — voir le membre polonais
%  de CENO et de SERIO.
%
%  LES MOIS ET LES JOURS S'ECRIVENT EN MINUSCULE en polonais, comme en
%  ido et en francais, et contrairement a l'anglais ou a l'allemand :
%  marzec, kwiecień, maj ; poniedziałek, wtorek, środa. Le gras du
%  releve les prend tels quels.
%
%  ET LE GENITIF POLONAIS RETOURNE CE QUE L'IDO POSE : « la flecho dil
%  klosheyo » donne « iglica dzwonnicy », tete d'abord. Les renvois
%  (7) et (5) restent donc dans l'ordre de l'ido ; la ou ils ne le
%  resteraient pas, la phrase se refait autour — jamais le renvoi.
% ===================================================================

%%K t03-apar-1 apar
\VUsaut{3.0mm}
\VUtitre{15.0pt}{60}{TABLICA Nr 3}
\VUsaut{1.6mm}
\VUtitre{11.4pt}{0}{Dzieciństwo i młodość.}
\VUsaut{3.4mm}

%%K t03-apar-2 apar
(Tablice 3 i 4 są tak zestawione, że przedstawiają w czterech \VUgras{scenach
rodzinnych} zasadnicze pojęcia o \VUgras{pogodzie}, \VUgras{wieku}, \VUgras{rodzinie},
\VUgras{ubiorze} i \VUgras{stanie}.)
\VUsaut{4.0mm}

%%K t03-tit-1 sub
\VUcentre{12.0pt}{0}{\textit{Scena pierwsza.}}
\VUsaut{3.0mm}

%%K t03-tit-2 sub
\VUpk{10.2pt}{40}{Uroczystość chrztu na wsi.}
\VUsaut{2.4mm}

%%K t03-tit-3 sub
\VUpk{10.2pt}{40}{Wiosna.}
\VUsaut{3.0mm}

%%K t03-c1-01-1 p
1. --- Jesteśmy na \VUgras{wiosnę}, w miesiącach \VUgras{marcu}, \VUgras{kwietniu} albo
\VUgras{maju}, w dniach \VUgras{tygodnia}: \VUgras{poniedziałek}, \VUgras{wtorek} albo
\VUgras{środa}. Jest \VUgras{dziesiąta} rano.

%%K t03-c1-02-1 p
2. --- \VUgras{Pogoda} jest piękna, \VUgras{powietrze} czyste. Słońce świeci. Kilka
\VUgras{obłoczków} \textsuperscript{(2)} wznosi się na niebieskim \VUgras{niebie}
\textsuperscript{(1)}.

%%K t03-c1-03-1 p
3. --- \VUgras{Drzewa} ledwie mają \VUgras{liście}. Smukłe \VUgras{topole}
\textsuperscript{(3)} i wysoki \VUgras{platan} \textsuperscript{(17)} mają dotąd tylko
pąki.

%%K t03-c1-04-1 p
4. --- \VUgras{Jaskółki} \textsuperscript{(4)} wracają i latają z radosnym świergotem
wokół \VUgras{iglicy} \textsuperscript{(7)} \VUgras{dzwonnicy} \textsuperscript{(5)},
która wieńczy \VUgras{kościół} \textsuperscript{(6)} wsi.

%%K t03-tit-4 sub
\VUsaut{3.4mm}
\VUpk{10.2pt}{40}{Kościół.}
\VUsaut{3.0mm}

%%K t03-c1-05-1 p
5. --- Bardzo skromny jest ten kościół ze swoim wielkim \VUgras{portalem}
\textsuperscript{(13)} i grubym \VUgras{zegarem} \textsuperscript{(8)}. Mała
\VUgras{wskazówka} \textsuperscript{(9)} jest na cyfrze X, wskazuje \VUgras{godziny}.
\VUgras{Duża} \textsuperscript{(10)} jest na XII (w połowie); wskazuje \VUgras{minuty}.

%%K t03-c1-06-1 p
6. --- W dzwonnicy \VUgras{dzwony} \textsuperscript{(11)} wesoło dzwonią. Na szczycie
dachu stoi mały kamienny \VUgras{krzyż} \textsuperscript{(12)}.

%%K t03-c1-07-1 p
7. --- Kościół ma nie tylko wielki \VUgras{portal} \textsuperscript{(13)}, ma też małe
drzwi boczne. Przez te drzwiczki ksiądz \VUgras{proboszcz} \textsuperscript{(15)},
ubrany w swoją \VUgras{sutannę}, wychodzi z \VUgras{zakrystii} \textsuperscript{(14)} i
idzie na \VUgras{plebanię} \textsuperscript{(19)}.

%%K t03-c1-08-1 p
8. --- Obok niego oto \VUgras{mleczarka} \textsuperscript{(24)} ze swoim
\VUgras{osłem} \textsuperscript{(23)}. Wraca z miasta, do którego zaniosła dobre mleko
dla dzieci.

%%K t03-tit-5 sub
\VUsaut{4.0mm}
\VUpk{10.2pt}{40}{Sad.}
\VUsaut{3.4mm}

%%K t03-c1-09-1 p
9. --- Pan Kłapouchy (osioł) jest całkiem zadowolony z tego porannego biegu. Z rozkoszą
wdycha powietrze przesycone wonią \VUgras{kwiatów} \textsuperscript{(20)}, które
okrywają \VUgras{drzewa owocowe} \textsuperscript{(18)} sąsiedniego \VUgras{sadu}.
Całkiem różowe i białe są te \VUgras{brzoskwinie} \textsuperscript{(18)} i
\VUgras{morele} \textsuperscript{(19)}, i budzą nadzieję dobrych zbiorów.

%%K t03-c1-10-1 p
10. --- Tymczasem małe \VUgras{pszczoły} \textsuperscript{(22)} z \VUgras{ula}
\textsuperscript{(21)} lecą tam brzęcząc na łup po swój słodki \VUgras{miód}.

%%K t03-c1-11-1 p
11. --- Ale co się dzieje? Oto dzieci ze wsi nadbiegają i płoszą przestraszone
\VUgras{gęsi} \textsuperscript{(25)}. To orszak wychodzi z kościoła po
\VUgras{chrzcie} syna notariusza.

%%K t03-c1-12-1 p
12. --- Mały Jakub w swojej \VUgras{kołysce} \textsuperscript{(26)} budzi się od
radosnego bicia dzwonów, a jego siostra Magdalena, która też chce zobaczyć orszak
chrzestny, prosi matkę, żeby przyszła uczesać jej włosy i ubrać ją na progu domu.

%%K t03-c1-13-1 p
13. --- Matka się zgadza. Wkłada swoją \VUgras{chustkę} \textsuperscript{(28)}, swój
\VUgras{stanik} \textsuperscript{(29)} i swój \VUgras{fartuch} \textsuperscript{(30)}, i
przychodzi usiąść na małym stołku przed drzwiami. Magdalena ma tylko swoją
\VUgras{halkę} \textsuperscript{(31)} i swój \VUgras{gorsecik} \textsuperscript{(32)}.

%%K t03-c1-14-1 p
14. --- Jakże jest szczęśliwa! Zapomina o swoich ulubionych kwiatach, o swoich
\VUgras{goździkach} \textsuperscript{(34)}, na których siadają \VUgras{motyle}
\textsuperscript{(35)}, o swoich \VUgras{tulipanach} \textsuperscript{(36)}, o swoich
\VUgras{konwaliach} \textsuperscript{(37)}, w \VUgras{doniczkach} \textsuperscript{(38)},
na \VUgras{murze} \textsuperscript{(33)}, i o swoich \VUgras{różach}
\textsuperscript{(39)}, które tworzą tak piękny żywopłot. Przygląda się bogatym strojom
pań z orszaku.

%%K t03-c1-15-1 p
15. --- Dzieci ze wsi niewiele zważają na stroje. Rzucają się i popychają jedno drugie,
żeby dostać \VUgras{cukierki} \textsuperscript{(55)} albo \VUgras{monety}
\textsuperscript{(52)}. Jedno z nich płacze \textsuperscript{(40)}.

%%K t03-c1-16-1 p
16. --- Drugie nadepnęło na łapę \VUgras{psa} \textsuperscript{(42)}, który ucieka
skowycząc i płoszy \VUgras{kurę} \textsuperscript{(43)} i jej \VUgras{kurczęta}
\textsuperscript{(44)}.

%%K t03-tit-6 sub
\VUsaut{5.0mm}
\VUpk{10.2pt}{40}{Orszak chrzestny.}
\VUsaut{4.0mm}

%%K t03-c1-17-1 p
17. --- Przed orszakiem idzie \VUgras{mamka} \textsuperscript{(45)}. Ma piękny
\VUgras{czepek} \textsuperscript{(46)} z długimi czerwonymi wstążkami. Niesie
\VUgras{noworodka} \textsuperscript{(47)} całego ubranego w \VUgras{koronki}
\textsuperscript{(48)}.

%%K t03-c1-18-1 p
18. --- Za mamką idą \VUgras{ojciec chrzestny} \textsuperscript{(49)} i \VUgras{matka
chrzestna} \textsuperscript{(53)}. Rozdają cukierki i monety. Ojciec chrzestny ma
\VUgras{rękawiczki} \textsuperscript{(51)} i \VUgras{cylinder} \textsuperscript{(50)}.
Matka chrzestna trzyma w ręku piękny \VUgras{bukiet} \textsuperscript{(54)}.

%%K t03-c1-19-1 p
19. --- Za nimi widzimy \VUgras{wuja} \textsuperscript{(56)} i \VUgras{ciotkę}
\textsuperscript{(58)} dziecka. Ciotka ma elegancki \VUgras{kapelusz}
\textsuperscript{(59)}, wuj czarny \VUgras{surdut} \textsuperscript{(57)} i białą
kamizelkę. Oto dalej \VUgras{kuzyn} \textsuperscript{(60)} i \VUgras{kuzynka}
\textsuperscript{(61)}. Ta trzyma za rękę \VUgras{siostrę} \textsuperscript{(62)}
noworodka, małą Bertę.

%%K t03-c1-20-1 p
20. --- Drugi \VUgras{brat} \textsuperscript{(63)} Berty idzie z tyłu. Podaje rękę
\VUgras{krewnej} \textsuperscript{(64)}. \VUgras{Przyjaciele} \textsuperscript{(65)}
rodziny idą potem.

%%K t03-tit-7 sub
\VUsaut{4.6mm}
\VUpk{10.2pt}{40}{Widzowie.}
\VUsaut{3.4mm}

%%K t03-c1-21-1 p
21. --- Obok orszaku oto \VUgras{żebrak} \textsuperscript{(66)} wsparty na swojej
\VUgras{kuli} \textsuperscript{(67)}. Prosi o jałmużnę (żebrze).

%%K t03-c1-22-1 p
22. --- Z drugiej strony jest mały chłopiec bardzo \VUgras{grzeczny}
\textsuperscript{(68)}. Kłania się i mówi « \VUgras{dzień dobry} » osobom, które zna.

%%K t03-c1-23-1 p
23. --- Mieszkańcy wsi wyszli ze swoich domów, żeby zobaczyć orszak chrzestny. Widzimy
\VUgras{wieśniaka} \textsuperscript{(69)} w swojej \VUgras{bawełnianej czapce}, a obok
\VUgras{wieśniaczkę} \textsuperscript{(70)}. Potem \VUgras{staruszkę}
\textsuperscript{(71)} wspartą na swojej \VUgras{lasce} \textsuperscript{(72)}. Wreszcie
\VUgras{młynarza} \textsuperscript{(73)} w swoim wielkim \VUgras{filcowym kapeluszu}
\textsuperscript{(74)} i młodą wieśniaczkę obutą w \VUgras{drewniaki}
\textsuperscript{(75)}, która uczy chodzić swoje dzieciątko.

%%K t03-c1-24-1 p
24. --- Ta scena jest bardzo zabawna.
\VUsaut{4.0mm}

%%K t03-tit-8 sub
\VUcentre{12.0pt}{0}{\textit{Scena druga.}}
\VUsaut{3.0mm}

%%K t03-tit-9 sub
\VUpk{10.2pt}{40}{Święto publiczne.}
\VUsaut{2.4mm}

%%K t03-tit-10 sub
\VUpk{10.2pt}{40}{Gry wodne i regaty.}
\VUsaut{3.0mm}

%%K t03-c2-01-1 p
1. --- Nadeszło \VUgras{lato}. Upalne słońce miesiąca \VUgras{czerwca} (lipca albo
\VUgras{sierpnia}) zachęca młodych do kąpieli i wioślarstwa.

%%K t03-c2-02-1 p
2. --- Na prawym brzegu rzeki wioślarze rozbierają się, żeby wziąć udział w grach
wodnych i regatach.

%%K t03-c2-03-1 p
3. --- Jeden z nich zdejmuje swoje \VUgras{buty} \textsuperscript{(1)}; rozwiązuje je
(rozpina). Jego dwaj koledzy są już gotowi. Odtąd mają tylko swoją \VUgras{trykotową
koszulkę} \textsuperscript{(2)} i \VUgras{kąpielówki} \textsuperscript{(3)}. Rozmawiają
czekając na swoich przyjaciół. Mówią o jarmarku i o święcie, które odbywa się na drugim
brzegu.

%%K t03-c2-04-1 p
4. --- Na ziemi, obok nich, leżą \VUgras{ubrania} ich kolegów: para \VUgras{trzewików}
\textsuperscript{(4)}, \VUgras{butów} \textsuperscript{(5)}, \VUgras{skarpetek}
\textsuperscript{(6)}, \VUgras{filcowe} \textsuperscript{(7)} i \VUgras{słomkowe}
\textsuperscript{(8)} kapelusze i \VUgras{krawat} \textsuperscript{(9)}.

%%K t03-c2-05-1 p
5. --- Jeden z młodzieńców starannie złożył swoją \VUgras{marynarkę}
\textsuperscript{(10)}. Kładzie ją na ręczniku, w którym jego sąsiad wkrótce zawinie
oba \VUgras{ubrania}.

%%K t03-c2-06-1 p
6. --- Szósty chłopiec się spóźnia. Ma jeszcze swoją \VUgras{czapkę}
\textsuperscript{(11)} na głowie. Nie zdjął ani swojej \VUgras{koszuli}
\textsuperscript{(12)}, ani swojego \VUgras{kołnierzyka} \textsuperscript{(13)}, ani
swoich \VUgras{mankietów} \textsuperscript{(14)}. Trzyma w ręku swoją
\VUgras{kamizelkę} \textsuperscript{(17)} i ma jeszcze swoje \VUgras{spodnie}
\textsuperscript{(16)} i \VUgras{szelki} \textsuperscript{(15)}. Na pewno nie będzie
gotowy na najbliższy wyścig.

%%K t03-c2-07-1 p
7. --- Obok młodzieńców ścieżka, po której bokach jest wiele \VUgras{ostów}
\textsuperscript{(18)}, prowadzi na brzeg rzeki, gdzie ojciec Jakub przygotowuje wśród
\VUgras{trzcin} \textsuperscript{(19)} \VUgras{ognie sztuczne} \textsuperscript{(20)},
które wystrzelą wieczorem.

%%K t03-tit-11 sub
\VUsaut{4.4mm}
\VUpk{10.2pt}{40}{Wyścig.}
\VUsaut{3.4mm}

%%K t03-c2-08-1 p
8. --- Na rzece kilka pań, osłaniając się swoimi parasolkami, przejeżdża się
\VUgras{łodzią} \textsuperscript{(21)}. Śledzą wyścig, który jest bardzo ciekawy. W
każdej \VUgras{jolce} \textsuperscript{(22)} czterech \VUgras{wioślarzy}
\textsuperscript{(24)} wiosłuje ze wszystkich sił, podczas gdy \VUgras{sternik}
\textsuperscript{(26)} trzyma \VUgras{ster} \textsuperscript{(25)} i kieruje kruchą
łodzią wioślarską. \VUgras{Wiosła} \textsuperscript{(23)} rytmicznie uderzają w wodę i
lekkie łodzie płyną z zawrotną szybkością.

%%K t03-c2-09-1 p
9. --- Kilku młodzieńców walczy obok filaru mostu o nagrodę \VUgras{masztu
bukszprytowego} \textsuperscript{(28)}. Jeden z nich ostrożnie idzie po giętkim i
śliskim maszcie, żeby dosięgnąć małej \VUgras{chorągiewki} \textsuperscript{(29)}
umocowanej na końcu. Jego nieszczęśliwy \VUgras{współzawodnik} \textsuperscript{(30)}
wpadł do wody. Płynie, żeby wrócić na brzeg.

%%K t03-c2-10-1 p
10. --- Starsi młodzieńcy \VUgras{urządzają turniej łodzi} \textsuperscript{(32)} obok
\VUgras{nabrzeża} \textsuperscript{(33)}. Wszystkim tym zabawom przygląda się liczna
\VUgras{publiczność} \textsuperscript{(31)} wsparta o \VUgras{poręcz} \VUgras{mostu}
\textsuperscript{(27)} i stłoczona na \VUgras{brzegu}. Kilku widzów odwraca się jednak,
żeby śledzić zabawę przy \VUgras{maszcie z nagrodą} \textsuperscript{(45)} albo żeby
podziwiać \VUgras{orszak burmistrza}, który idzie na \VUgras{rozdanie} nagród.

%%K t03-c2-11-1 p
11. --- Najpierw oto kilku \VUgras{urwisów} \textsuperscript{(35)} idących przed
\VUgras{muzykami} \textsuperscript{(36)}; potem idzie \VUgras{burmistrz}
\textsuperscript{(37)}, jego zastępcy i \VUgras{rada gminna} miasteczka. Na końcu idą
\VUgras{strażacy} \textsuperscript{(38)}.

%%K t03-tit-12 sub
\VUsaut{5.0mm}
\VUpk{10.2pt}{40}{Jarmark.}
\VUsaut{3.6mm}

%%K t03-c2-12-1 p
12. --- Wielu ludzi jest na \VUgras{placu jarmarcznym} \textsuperscript{(39)}. Przychodzi
się zobaczyć rozmaite \VUgras{budy}: \VUgras{drewniane konie} \textsuperscript{(40)},
\VUgras{cyrk} \textsuperscript{(41)}, w którym przedstawienie już się zaczęło;
\VUgras{bal publiczny} \textsuperscript{(42)}, gdzie młodzieńcy i panny z miasta
zażywają przyjemności \VUgras{tańca}.

%%K t03-c2-13-1 p
13. --- W głębi widzimy \VUgras{arenę} \textsuperscript{(44)}, gdzie dzielny hiszpański
\VUgras{matador} walczy z rozjuszonym \VUgras{bykiem} \textsuperscript{(45)}, potem
\VUgras{kolejkę górską} \textsuperscript{(49)}, \VUgras{strzelnicę}
\textsuperscript{(46)}, gdzie ćwiczy się w używaniu \VUgras{broni palnej} i w strzelaniu
do \VUgras{celu}.

%%K t03-c2-14-1 p
14. --- Obok oto \VUgras{teatr jarmarczny} \textsuperscript{(47)}, w którym grywa się
\VUgras{komedię} i \VUgras{dramat}. Całkiem blisko jest buda \VUgras{olbrzyma}
\textsuperscript{(48)} i \VUgras{karła}. \VUgras{Wzrost} pierwszego wynosi dwa metry i
piętnaście centymetrów; drugi jest ledwie tak duży jak dziecko.

%%K t03-c2-15-1 p
15. --- Wreszcie oto wielka \VUgras{menażeria} \textsuperscript{(50)} ze swoimi
\VUgras{dzikimi zwierzętami}, swoim \VUgras{tygrysem} \textsuperscript{(51)} o
potężnych \VUgras{pazurach}, swoim \VUgras{lwem} \textsuperscript{(52)} o gęstej
\VUgras{grzywie}, swoimi dwiema \VUgras{żyrafami} \textsuperscript{(53)} i swoim
\VUgras{słoniem} \textsuperscript{(54)}, który tak zręcznie używa swojej
\VUgras{trąby}.

%%K t03-c2-16-1 p
16. --- \VUgras{Pogromca} \textsuperscript{(55)}, który tresuje te zwierzęta, stoi przy
drzwiach budy.
\VUsaut{6.0mm}
\VUfilet{22mm}
